%=====================================================================
\chapter{Normal Forms: 1NF to BCNF}\label{ch:normalforms}
%=====================================================================
\locator{5}{Part V --- Getting the Design Right}

\begin{outcomes}
By the end of this chapter you will be able to:
\begin{tight}
  \item \textbf{Name} the three anomalies --- update, insertion, deletion ---
        and \textbf{show} each one in a real table.
  \item \textbf{State} the conditions for 1NF, 2NF, 3NF and BCNF, and
        \textbf{test} a relation against each using its dependencies.
  \item \textbf{Identify} a partial and a transitive dependency, and
        \textbf{say} which normal form each one violates.
  \item \textbf{Decompose} a relation into BCNF by splitting on a violating
        dependency.
  \item \textbf{Weigh} what normalisation costs in reads against what it saves
        in correctness, using measured figures.
\end{tight}
\end{outcomes}

\begin{prereq}
\chref{fds}: functional dependencies, attribute closure and candidate keys ---
every test in this chapter is a closure. The \textit{report} table from that
chapter's laboratory. You need \texttt{code/m18\_normalforms.py}.
\end{prereq}

\begin{keyterms}
normalisation $\bullet$ anomaly $\bullet$ update anomaly $\bullet$ insertion
anomaly $\bullet$ deletion anomaly $\bullet$ first normal form $\bullet$
atomic $\bullet$ partial dependency $\bullet$ second normal form $\bullet$
transitive dependency $\bullet$ third normal form $\bullet$ prime attribute
$\bullet$ Boyce--Codd normal form $\bullet$ decomposition $\bullet$ redundancy
\end{keyterms}

\section{Three Ways a Wide Table Goes Wrong}

\chref{fds} built \textit{report}: fourteen columns, one row per enrolment,
every fact the university holds about it in one place. It answers questions
without joins. This chapter is about what it does when the data
\emph{changes}, which is the thing it was never designed for.

\begin{definitionbox}[The three anomalies]
An \term{update anomaly}: one fact is stored in many rows, so changing it means
changing all of them --- and changing some of them leaves the table holding two
versions of the truth.

\smallskip
An \term{insertion anomaly}: a fact cannot be recorded without also recording
an unrelated one that does not exist yet.

\smallskip
A \term{deletion anomaly}: deleting one fact removes another that had nowhere
else to live.

\smallskip
All three have one cause. \textbf{A fact whose determinant is smaller than the
key is stored once per key value}, so the table is holding several relations
at once and can only change them together.
\end{definitionbox}

\begin{measured}[The same information, held wide and held normalised]
\textsf{\footnotesize \texttt{code/m18\_normalforms.py}, seeded database
(4{,}000 students), SQLite 3.51.2, Intel i5-6600, Windows 11. Best of three
runs where timed.}

\medskip
\textbf{1. Bytes, and copies of one fact}

\begin{tabular}{@{}L{6.6cm}R{2.0cm}R{2.6cm}L{2.8cm}@{}}
& \textbf{rows} & \textbf{bytes} & \\
\textit{report}, as \sqlkw{CREATE TABLE AS} made it & 60{,}037 & 6{,}860{,}800 & no index \\
\textit{report}, with its key indexed & 60{,}037 & 9{,}289{,}728 & like for like \\
normalised, 6 tables & 65{,}693 & 5{,}574{,}656 & key indexes \\
\end{tabular}

\smallskip
\textbf{Like for like, the wide table is 1.7$\times$ the size.} Compared as
created it looks like 1.2$\times$, because \sqlkw{CREATE TABLE AS} copies rows
and nothing else --- no key, no index --- while every normalised table carries
the index its \sqlkw{PRIMARY KEY} implies. The bytes are not the real cost;
the copies are:

\begin{tabular}{@{}L{6.6cm}R{2.6cm}R{2.6cm}@{}}
\textbf{copies of one fact} & \textbf{\textit{report}} & \textbf{normalised} \\
a student's name & 18 & 1 \\
an instructor's name & 949 & 1 \\
a course's title & 1{,}382 & 1 \\
a department's name & \textbf{7{,}936} & 1 \\
\end{tabular}

\medskip
\textbf{2. The update anomaly} --- rename one department

\begin{tabular}{@{}L{6.6cm}R{2.6cm}R{2.6cm}@{}}
& \textbf{rows changed} & \textbf{ms} \\
\textit{report} & 7{,}869 & 378.9 \\
\textit{department} & 1 & 25.4 \\
\end{tabular}

\smallskip
Most of the 25.4\,ms is the commit, not the change: one row costs about as
much to make durable as it does to write.

\medskip
\textbf{3. A half-finished update} --- the IT department renamed in the
FALL2025 rows only, as a crashed program or a wrong \sqlkw{WHERE} would leave
it

\begin{tabular}{@{}L{6.6cm}R{2.6cm}@{}}
\texttt{'Information Technology'} & 5{,}273 rows \\
\texttt{'Information Tech.'} & 2{,}663 rows \\
\end{tabular}

\smallskip
Two names for one department, in one table, and no constraint anywhere that
could have refused it. In \textit{department} the same mistake is impossible:
there is only one place to write the name.

\medskip
\textbf{4. Insertion and deletion}

\begin{tabular}{@{}L{8.0cm}L{6.0cm}@{}}
a new course, nobody enrolled yet: rows in \textit{report} that can hold it &
  \textbf{0} \\
the same course in \textit{course} & one row, accepted \\
\textit{report} after deleting the 1{,}122 enrolments of PH-350: rows still
  describing the course & \textbf{0} \\
the same course in \textit{course} & still there \\
\end{tabular}

\medskip
\textbf{5. Read cost} --- where the wide table is supposed to win

\begin{tabular}{@{}L{8.0cm}R{2.4cm}L{3.6cm}@{}}
\textbf{(a) one student's transcript} & \textbf{ms} & \\
\textit{report}, no index on roll\_no & 15.35 & \textbf{162$\times$ slower} \\
normalised, one join & 0.09 & \\
\textit{report}, with an index on roll\_no & 0.08 & a tie \\
\textbf{(b) enrolments per department} & & \\
\textit{report}, no joins & 66.51 & \textbf{1.5$\times$ faster} \\
normalised, two joins & 98.34 & \\
\end{tabular}

\smallskip
\textbf{The wide table lost the lookup by two orders of magnitude}, and not
because of a join: the normalised \textit{enrolment} table's primary key
begins with \texttt{roll\_no}, so the transcript was an index search, while
\textit{report} had no index and was read end to end. Given one index the two
tied. It won only the full-scan report, by 1.5$\times$ --- the case where
reading one table instead of joining three is a real saving.
\end{measured}

\begin{alertbox}[So the trade is not ``fast reads against slow writes'']
That is how it is usually described, and the measurement does not support it:

\begin{tight}
  \item \textbf{Lookups were not faster wide.} They were 162$\times$ slower
        until an index was added, and equal after. ``No joins'' is not an
        advantage when the join is an index search.
  \item \textbf{Full scans were faster wide, by 1.5$\times$.} Real, and modest.
  \item \textbf{Writes were 7{,}869 times the work}, and one half-finished
        write left the table permanently self-contradictory.
\end{tight}

\smallskip
The honest summary: normalisation costs a little on reports over every row,
and buys correctness that no amount of care in the application can supply.
\chref{decomposition} asks when the 1.5$\times$ is worth paying for, and the
answer is a measurement, not a preference.
\end{alertbox}

\section{The Normal Forms}

Each normal form forbids one kind of dependency. They are nested: every BCNF
relation is in 3NF, every 3NF relation in 2NF, and every 2NF relation in 1NF.

\begin{definitionbox}[1NF and 2NF]
\textbf{First normal form (1NF):} every attribute value is atomic --- no
lists, no repeating groups, no ``three phone numbers in one field''. It is
\chref{relational}'s fourth property of a relation, and \chref{mapping}'s rule
for multi-valued attributes.

\smallskip
A \term{partial dependency} is one whose determinant is a \emph{proper subset}
of a candidate key: \fd{roll\_no}{name}, when the key is \{roll\_no,
course\_code, semester, section\_no\}.

\smallskip
\textbf{Second normal form (2NF):} 1NF, and no non-prime attribute is partially
dependent on any candidate key. (A \term{prime attribute} is one that belongs
to some candidate key.)
\end{definitionbox}

\begin{definitionbox}[3NF and BCNF]
A \term{transitive dependency} is \fd{K}{Y} holding only because \fd{K}{X} and
\fd{X}{Y}, where $X$ is not a superkey: \fd{course\_code}{dept\_name} through
\textit{dept\_code}.

\smallskip
\textbf{Third normal form (3NF):} for every non-trivial \fd{X}{A}, either $X$
is a superkey, or $A$ is a prime attribute.

\smallskip
\textbf{Boyce--Codd normal form (BCNF):} for every non-trivial \fd{X}{A},
$X$ is a superkey. Full stop --- the prime-attribute exception is gone.

\smallskip
\textbf{The one sentence to remember:} in BCNF, every determinant is a
superkey. Every fact is about the key, the whole key, and nothing but the key.
\end{definitionbox}

\begin{worked}[\textit{report}, from nowhere to BCNF]
Key $K = \{$roll\_no, course\_code, semester, section\_no$\}$ (\chref{fds}).
Test each dependency: is its determinant a superkey?

\medskip
\begin{tabular}{@{}L{6.2cm}L{2.5cm}L{5.3cm}@{}}
\textbf{dependency} & \textbf{superkey?} & \textbf{violates} \\
\fd{roll\_no}{name, programme} & no & 2NF --- partial \\
\fd{course\_code}{title, credit\_hours, dept\_code} & no & 2NF --- partial \\
\fd{dept\_code}{dept\_name} & no & 3NF --- transitive \\
\fd{course\_code, semester, section\_no}{room, emp\_id} & no & 2NF --- partial \\
\fd{emp\_id}{instructor} & no & 3NF --- transitive \\
\fd{$K$}{grade} & \textbf{yes} & nothing \\
\end{tabular}

\medskip
Five violations, one clean dependency. \textbf{The decomposition procedure:}
take a violating \fd{X}{Y}; split $R$ into $R_1 = X \cup Y$ and $R_2 = R - Y$;
repeat on whatever still violates. \textbf{Split on the transitive
dependencies first} --- the box below shows what goes wrong otherwise.

\medskip
\begin{tightnum}
  \item On \fd{dept\_code}{dept\_name}: \relation{department}{\relkey{dept\_code},
        dept\_name}. Remove dept\_name from the rest.
  \item On \fd{emp\_id}{instructor}: \relation{instructor}{\relkey{emp\_id},
        instructor}. Remove instructor.
  \item On \fd{roll\_no}{name, programme}: \relation{student}{\relkey{roll\_no},
        name, programme}. Remove name and programme.
  \item On \fd{course\_code}{title, credit\_hours, dept\_code}:
        \relation{course}{\relkey{course\_code}, title, credit\_hours,
        dept\_code}. Remove title, credit\_hours and dept\_code.
  \item On \fd{course\_code, semester, section\_no}{room, emp\_id}:
        \relation{section}{\relkey{course\_code}, \relkey{semester},
        \relkey{section\_no}, room, emp\_id}. Remove room and emp\_id.
  \item What remains: \relation{enrolment}{\relkey{roll\_no},
        \relkey{course\_code}, \relkey{semester}, \relkey{section\_no},
        grade} --- whose only dependency has the key as its determinant.
\end{tightnum}

\medskip
\textbf{Why the order matters.} Split on \fd{course\_code}{title, credit\_hours,
dept\_code} \emph{first} and $R_2 = R - Y$ removes \textit{dept\_code} from the
remainder --- but leaves \textit{dept\_name} there, with nothing beside it that
determines it except, transitively, \textit{course\_code}. The next split is
then on \fd{course\_code}{dept\_name}, and produces
\relation{}{\relkey{course\_code}, dept\_name}: a department's name stored
once \emph{per course}, which is the redundancy you were removing. Every
relation is still in BCNF, so no normal-form test fails --- but the department
has no table of its own, and \fd{dept\_code}{dept\_name} now lives in
\emph{no} single table, so nothing can enforce it: two CS courses could carry
two different department names. \chref{decomposition} calls that losing
dependency preservation. Handle the deepest dependency first and each fact
lands with its own determinant.

\medskip
\textbf{Six relations, every one in BCNF --- and they are exactly the six
tables \chref{mapping} produced from the ER diagram.} Two routes, one from
design and one from dependencies, arrived at the same schema. That agreement
is the best evidence a design is right.
\end{worked}

\dsfig{%
\begin{tikzpicture}[
  wide/.style={rounded corners=3pt, draw=accent, line width=1.1pt,
               fill=accent!7, text=accent!70!black, font=\scriptsize,
               align=left, text width=4.4cm, inner sep=4pt},
  tb/.style={rounded corners=2pt, draw=#1, line width=0.9pt, fill=#1!7,
             text=#1!60!black, font=\scriptsize, align=left,
             text width=4.5cm, inner sep=3pt},
  arr/.style={-{Stealth[length=2.2mm]}, primary!50, line width=0.9pt},
  lbl/.style={font=\scriptsize\itshape, text=neutral, align=center}
]
  \node[wide] (r) at (-4.9,0) {\textbf{report} --- 14 attributes\\[2pt]
       roll\_no, name, programme, course\_code, title, credit\_hours,
       dept\_code, dept\_name, semester, section\_no, room, emp\_id,
       instructor, grade\\[3pt]
       \textbf{60{,}037 rows}; a department's name stored
       \textbf{7{,}936} times};

  \node[tb=primary]  (s) at (2.6, 2.25) {\textbf{student}(\relkey{roll\_no},
       name, programme)};
  \node[tb=primary]  (c) at (2.6, 1.35) {\textbf{course}(\relkey{course\_code},
       title, credit\_hours, dept\_code)};
  \node[tb=greenhl]  (d) at (2.6, 0.45) {\textbf{department}(\relkey{dept\_code},
       dept\_name)};
  \node[tb=purplehl] (x) at (2.6,-0.45) {\textbf{section}(\relkey{course\_code},
       \relkey{semester}, \relkey{section\_no}, room, emp\_id)};
  \node[tb=greenhl]  (i) at (2.6,-1.35) {\textbf{instructor}(\relkey{emp\_id},
       instructor)};
  \node[tb=accent]   (e) at (2.6,-2.25) {\textbf{enrolment}(\relkey{roll\_no},
       \relkey{course\_code}, \relkey{semester}, \relkey{section\_no},
       grade)};

  \foreach \t in {s,c,x,e}{\draw[arr] (r.east) -- (\t.west);}
  \draw[arr, greenhl!60] (c.east) .. controls (5.4,1.35) and (5.4,0.45)
       .. (d.east);
  \draw[arr, greenhl!60] (x.east) .. controls (5.4,-0.45) and (5.4,-1.35)
       .. (i.east);

  \node[lbl, text=primary!70!black] at (-1.2,2.35) {partial: 2NF};
  \node[lbl, text=greenhl!60!black, anchor=west] at (5.6,0.9)
       {transitive:\\3NF};
  \node[lbl, text=greenhl!60!black, anchor=west] at (5.6,-0.9)
       {transitive:\\3NF};
\end{tikzpicture}}%
{\textit{report} decomposed to BCNF. Splitting on each partial dependency
gives four relations; splitting on the two transitive ones gives two more. The
result is the six tables \chref{mapping} derived from the ER diagram by an
entirely different route.}{bcnf}

\begin{conceptbox}[3NF against BCNF: the one case where they differ]
They differ only when a relation has two \emph{overlapping} composite
candidate keys --- rare, but it exists. A classic case:
\relation{teaches}{student, course, instructor}, where each instructor teaches
one course (\fd{instructor}{course}) and each student has one instructor per
course (\fd{student, course}{instructor}).

\smallskip
\fd{instructor}{course}: \textit{instructor} is not a superkey, so BCNF is
violated. But \textit{course} is prime --- it is in the key \{student,
course\} --- so 3NF is satisfied. Decomposing to BCNF gives
\relation{}{\relkey{instructor}, course} and \relation{}{\relkey{student},
\relkey{instructor}}, and \textbf{loses} the dependency
\fd{student, course}{instructor}: no single table can now enforce it.

\smallskip
That trade --- BCNF, or keeping every dependency enforceable --- is
\chref{decomposition}'s subject. For the university schema it never arises.
\end{conceptbox}

\begin{pitfall}
\begin{tight}
  \item \textbf{Thinking normalisation is about saving bytes.} It saved a
        factor of 1.7. It removed 7{,}935 copies of one department name.
        The copies are the point.
  \item \textbf{Believing a wide table reads faster.} Measured: 162$\times$
        slower on a lookup without an index; a tie with one.
  \item \textbf{Comparing a \sqlkw{CREATE TABLE AS} copy with indexed tables.}
        The copy has no index. Measure like for like.
  \item \textbf{Testing 2NF against one key when there are several.} A partial
        dependency on \emph{any} candidate key violates it.
  \item \textbf{Stopping at 3NF without checking BCNF.} They usually coincide;
        check rather than assume.
  \item \textbf{Splitting a relation and losing the join.} Every decomposition
        must keep the determinant in both halves --- \chref{decomposition}
        explains why.
\end{tight}
\end{pitfall}

\begin{lab}[Break the wide table three ways, then decompose it]
\begin{lstlisting}[language=bash]
# --- 1. Seed, and rebuild report from Chapter 17 --------------------
python code/schema/seed.py university.db
sqlite3 university.db
\end{lstlisting}
\begin{lstlisting}[language=sql]
-- --- 2. Recreate report (Chapter 17's lab, step 2), then ---------
SELECT COUNT(*) FROM report;

-- --- 3. The update anomaly -------------------------------------
.timer on
UPDATE report SET dept_name = 'Computing' WHERE dept_code = 'CS';
-- Rows changed? (.changes on, or SELECT changes();)
UPDATE department SET name = 'Computing' WHERE dept_code = 'CS';

-- --- 4. Leave it half done --------------------------------------
UPDATE report SET dept_name = 'Information Tech.'
 WHERE dept_code = 'IT' AND semester = 'FALL2025';
SELECT dept_name, COUNT(*) FROM report
 WHERE dept_code = 'IT' GROUP BY dept_name;
-- Two answers. Write the query a programmer would need to find out
-- which name is current. Can it be written at all?

-- --- 5. The insertion anomaly -----------------------------------
-- Add a course nobody is enrolled in yet, to report.
-- Try it. What would you have to put in roll_no and grade?

-- --- 6. The deletion anomaly ------------------------------------
SELECT course_code, COUNT(*) FROM report
 GROUP BY course_code ORDER BY 2 LIMIT 1;
-- Delete every row for that course. Then ask report its title.

-- --- 7. Read cost, like for like --------------------------------
SELECT course_code, title FROM report
 WHERE roll_no = (SELECT MIN(roll_no) FROM student);
CREATE INDEX idx_report_roll ON report(roll_no);
SELECT course_code, title FROM report
 WHERE roll_no = (SELECT MIN(roll_no) FROM student);
-- Two timings. Which part of the difference was the join?
\end{lstlisting}
\begin{lstlisting}[language=bash]
# --- 8. On paper: decompose report to BCNF ---------------------------
# One violating dependency at a time. Name each new relation, its key,
# and which normal form the split repaired. Compare your six relations
# with Figure 9.1. Then run the measured version:
python code/m18_normalforms.py
\end{lstlisting}

\smallskip
\textbf{What to notice.} Step 4 is the one with no repair. Once a table holds
two versions of a fact and nothing records which came later, no query can tell
you which is true. That is not a performance problem, and no index fixes it.
It is the reason normal forms exist.
\end{lab}

\begin{checkpoint}
\begin{tightnum}
  \item For each of the three anomalies, give the specific operation on
        \textit{report} that exhibits it, the measured consequence, and the
        single dependency responsible.
  \item \relation{R}{\relkey{A}, \relkey{B}, C, D} has \fd{A}{C} and
        \fd{C}{D}. State its normal form, identify each violation by kind, and
        decompose it to BCNF.
  \item The wide table read a transcript 162$\times$ slower than the
        normalised schema and a department report 1.5$\times$ faster. Explain
        both results, and say which one would change if \textit{report} were
        given an index on \texttt{roll\_no}, and to what.
\end{tightnum}
\end{checkpoint}

\begin{chaptersummary}
\begin{tight}
  \item Three anomalies --- \textbf{update}, \textbf{insertion},
        \textbf{deletion} --- share one cause: a fact whose determinant is
        smaller than the key is stored once per key value.
  \item \textbf{Measured:} one department's name was held \textbf{7{,}936}
        times; renaming it changed \textbf{7{,}869} rows against 1; a
        half-finished update left \textbf{two names} for one department.
  \item \textbf{Measured:} a new course could not be inserted into the wide
        table at all, and deleting a course's enrolments deleted the course.
  \item \textbf{Measured:} like for like, the wide table was
        \textbf{1.7$\times$} the size --- 1.2$\times$ only when compared
        without its index.
  \item \textbf{Measured:} the wide table was \textbf{162$\times$ slower} on a
        lookup with no index and tied with one; it was \textbf{1.5$\times$
        faster} only on a full-scan report.
  \item \textbf{1NF}: atomic values. \textbf{2NF}: no partial dependency.
        \textbf{3NF}: no transitive dependency on a non-prime attribute.
        \textbf{BCNF}: every determinant is a superkey.
  \item Decompose on a violating \fd{X}{Y} into $X \cup Y$ and $R - Y$, and
        repeat.
  \item \textit{report} decomposes to exactly the six tables \chref{mapping}
        derived from the ER diagram --- two independent routes, one schema.
  \item 3NF and BCNF differ only with overlapping composite keys, where BCNF
        can cost an enforceable dependency.
\end{tight}
\end{chaptersummary}

\begin{reviewq}
  \item \textbf{(MCQ)} A non-prime attribute depending on part of a composite
        key violates: (a)~1NF; (b)~2NF; (c)~3NF but not 2NF; (d)~BCNF only.
  \item \textbf{(MCQ)} In BCNF, for every non-trivial \fd{X}{A}: (a)~$A$ is
        prime; (b)~$X$ is a superkey; (c)~$X$ is a candidate key;
        (d)~$A$ is in $X$.
  \item \textbf{(Short)} Why can no query repair a table holding two versions
        of the same fact?
  \item \textbf{(Short)} Give the dependency in \textit{report} that is
        transitive, and say which normal form it violates.
  \item \textbf{(Short)} Why was the wide table 1.2$\times$ the size as
        created but 1.7$\times$ like for like?
  \item \textbf{(Applied)} A spreadsheet records
        \relation{loan}{member\_id, member\_name, book\_id, title, author,
        date\_out, date\_back}, one row per loan. Write its dependencies,
        identify its key, classify every violation, and decompose it to BCNF,
        naming each relation's key.
  \item \textbf{(Applied)} A manager proposes keeping \textit{report} as the
        only table ``because reports are faster''. Using every figure in the
        measured box, write a reply that concedes what is true, refutes what is
        not, and proposes a design that keeps the 1.5$\times$ for reports
        without keeping the anomalies --- a design \chref{decomposition}
        measures.
\end{reviewq}
